\thispagestyle{plain}

Prediction markets aggregate collective intelligence to forecast uncertain events, but their utility depends on reliable resolution of outcomes. Current oracle systems force a tradeoff between fast but brittle automation and accurate but costly human arbitration. Recent work has shown that single-LLM oracles can achieve meaningful accuracy, but these systems inherit all failure modes of their underlying model with no mechanism for self-correction. This thesis evaluates whether multi-agent LLM architectures can improve oracle resolution accuracy over single-model baselines. We compare two architectures, independent aggregation and deliberative consensus, against single-LLM baselines using three diverse models (GPT-5 Nano, DeepSeek V3, Llama-3.3-70B) on 1,189 resolved prediction market questions from KalshiBench. All agents share a common evidence layer through Exa, with retrieval filtered by publication date to isolate reasoning capability from retrieval quality. Independent aggregation with confidence-weighted voting achieves the highest accuracy at 83.43\%, outperforming the best individual model by 1.01 percentage points. Deliberative consensus degrades accuracy to approximately 76\%, below every single-model baseline, a result we attribute to persuasive error propagation, where confidently wrong models flip correct ones during debate. Moderate-to-high error correlations across models ($r = 0.529$--$0.689$) explain why aggregation gains fall well short of the theoretical Condorcet ceiling and place a fundamental limit on ensemble-based approaches. A substantial fraction of questions resist correction by any multi-agent architecture, motivating principled escalation to human arbitration. Based on these findings, we propose routing criteria for hybrid AI-human oracle systems: auto-resolving only unanimous, high-confidence questions yields 97.87\% accuracy on 47\% of the dataset, with inter-agent disagreement flagging the remainder for human arbitration.
